% front_a: 题词 + 前言 (原书 pVI 题词页 / 书页 vii–x / PDF p006, p008–p011)
% 本文件经 scripts/assemble.py 装配为 chapters/ch00_front.tex

% 原书题词页（单独一页）
\clearpage
\epigraphbook{因为，倘若每颗星星不过是数学上的一个点，\\
凭借赤经与赤纬定位于天空半球之上，\\
那么所有的星星，纵然不可胜数，\\
也必像任何其他的点集一样，\\
转而代表某个单一的巨型方程——\\
在上帝的心智里，它像，比方说，\\
球面的方程一样直白，——\\
于我们则不可读、不可算。\\
这是一桩孤独、没有回报、甚或不可能完成的任务，——\\
然而我猜想，总得有些人永远在探寻。}%
{Thomas Pynchon，《Mason \& Dixon》}
\clearpage

\chapter*{前言}
\markboth{前言}{}
\addcontentsline{toc}{chapter}{前言}

广义相对论是有史以来所发明的最美丽的物理理论。它把我们经验世界中最普遍的特征之一——引力——用一种优雅的数学结构——弯曲时空的微分几何——来描述，给出毫不含糊的预言，并得到了令人惊叹的实验证实。广义相对论的各种推论，从大爆炸到黑洞，常常是年轻人最初对物理学产生兴趣的原因；而当求学者终于到达能够在严格的定量水平上理解这些现象的节点时，那种纯粹的喜悦难以言表。如果你正考虑读这本书，这个节点就在此处。

近几十年来，广义相对论（GR）已成为现代物理学不可分割的组成部分。自 Einstein 于 1916 年提出这一理论之后的很长一段时间里，GR 都被视为一项辉煌的成就，却在某种程度上处于有趣研究的主流之外。然而如今，越来越多不同专业的学生发现有必要学习 Einstein 的理论。GR 不仅本身是一个活跃的研究领域，它还是所有对天体物理、宇宙学、弦论乃至粒子物理感兴趣的人的标准课程的一部分。这并非要轻视 GR 更为实用的用途，其中包括全球定位系统（GPS）卫星网络的工作原理。

关于 GR 的书并不少，其中许多都很出色。事实上，大约三十年前，这门学科一口气出现了不下三本书，每一本都已成为经典：Weinberg（1972）、Misner、Thorne 与 Wheeler（1973），以及 Hawking 与 Ellis（1975）。这些书都浸透着作者本人强烈坚持的观点。这使得这些著作与它们的读者之间形成了一种爱恨交织的关系；对其中每一本来说，都不难找到宣称它是“有史以来最好的教科书”的学生，也不难找到觉得它完全无法下咽的学生。对持这些看法的个人而言，他们的判断很可能都是对的；进入这门学科的途径有很多种。

本书只有一个目的：为广义相对论提供一个清晰的导论，适合研究生或高年级本科生使用。我试图收入足够多的材料，使得几乎任何一门一学期的 GR 导论课程都能在本书中找到相应的内容，但也不再多出太多。特别是，我一直克制着去写一本包罗万象的参考书的诱惑。本书唯一的目标就是教你 GR。

我有意努力取常规而舍标新立异。如果说我的思想倾向有什么可指摘之处，那就是倾向于把广义相对论当作一种场论来思考——这一视角有助于领会 GR、粒子物理与弦论之间的联系。与此同时，GR 还有许多令人兴奋的天体物理应用（黑洞、引力透镜、引力波的产生与探测、早期宇宙、晚期宇宙、宇宙学常数），我至少尽力收入了关于这些论题的足够背景讨论，使学生能够着手研读当前文献。

任何一部广义相对论导论所面临的首要问题，是数学严格性的水准。对此并没有唯一正确的方案，因为不同的学生对不同方法的理解与热情各有不同。有鉴于此，我努力做到让每个人都各有所得。我没有回避细致的形式化，但也试图收入具体例子，并对所讨论的概念作非正式的讨论。大部分最数学化的材料被移到了附录。附录中的某些内容其实是这门课程不可分割的一部分（例如关于共形图的讨论），但每位读者或讲授者可以自行决定何时深入其中；正文里设有指示路标。

令人惊讶的是，学习广义相对论所需的形式先修课程非常少；大部分材料都是随本书的展开而建立起来的。这里绝不假定读者以前学过黎曼几何，学了也未必有帮助。如果已经学过一点狭义相对论会更好；尽管第 1 章中有相应的讨论，其目的更多在于回顾基础并引入若干记号，而不是提供一个自足的导论。除此之外，接触过一些电磁学、拉格朗日力学与线性代数或许有用，但其要点已包含在本书之中。

本书的结构应当是清楚的。第 1 章回顾狭义相对论与基本张量代数，包括对经典场论的简要讨论。接下来两章较为详细地引入流形与曲率；其中包含一些有启发性的物理，但首要目标是搭建数学框架。广义相对论本体在第 4 章引入，并顺带讨论了一些替代理论。随后四章讨论 GR 的三大应用：黑洞（两章）、微扰理论与引力波，以及宇宙学。这些论题近年来都经历了研究的大爆发，因此这里的讨论必然只是导引性的，但我尽量强调了与当前工作相关的议题。这三大应用可以按任何顺序讲授，尽管它们之间存在文中已标明的相互依赖。实验检验的讨论穿插于这些章节之中。第 9 章是弯曲时空量子场论的简要介绍；它并非初学 GR 的必要内容，但它对量子引力与宇宙学方面的工作已变得日益重要，因此值得一讲。另一方面，也有一些论题被“可耻地”忽略了：初值问题与宇宙学微扰论浮现在眼前，此外还有别的。幸而其他资源并不匮乏。附录各有其用途：既有正文中刻意回避的技术性讨论，也有本可放在多处的关键概念，还有一些有用但在主线之外的附加专题。

既然本书的目标是教学而非原创，当别的书（列于文献目录中）的讲法在我看来完全合理时，我常常大量倚重它们。倚重特别重的地方，我已在正文中标明。显然，一个主要资源是 Wald（1984）的书，它已成为本领域的标准参考；希望本书的读者能由此做好准备，得以进入 Wald 书中更高级的章节。

本书源自我在 MIT 讲授 GR 课程时准备的一组讲义。这些讲义在网上免费公开，并将继续如此；它们会链接到下面列出的网站。本书内容大概有一半多一点包含在讲义里，尽管拥有一本书（哪怕好几本）的好处应当是不言自明的。

无数人对我个人理解广义相对论、也对本书本身做出过重大贡献——多到无法指望一一致谢而求周全。但有几个人值得特别提及。Ted Pyne 与我一同学习这门学科，教了我许多，并在我们第一次讲授 GR 课程（哈佛大学天文系的一门研讨课）时与我合作；本书部分内容基于我们共同的笔记。Nick Warner 在 MIT 讲授了我最初学习 GR 的那门课，他的讲义无疑对本书的内容产生了极深的影响。Neil Cornish 好心提供了大量习题，其中许多已收在各章之末。在阅读过部分书稿并提出建议的众多人士中，Sanaz Arkani-Hamed 好心把整部书稿极为仔细地过了一遍。

我还要感谢所有当面或通过电子邮件对本书不同部分提出意见的人；其中包括 Tigran Aivazian、Teodora Beloreshka、Ed Bertschinger、Patrick Brady、Peter Brown、Jennifer Chen、Michele Ferraz Figueiró、Eanna Flanagan、Jacques Fric、Ygor Geurts、Marco Godina、Monica Guica、Jim Hartle、Tamás Hauer、Daniel Holz、Ted Jacobson、Akash Kansagra、Chuck Keeton、Arthur Kosowsky、Eugene Lim、Jorma Louko、Robert A.~McNees、Hayri Mutluay、Simon Ross、Itai Seggev、Robert Wald 与 Barton Zwiebach。若有遗漏未提及者，谨致歉意。在这一路上，我有幸从许多人那里获得智慧与见地，其中包括 Shadi Bartsch、George Field、Deryn Fogg、Ilana Harrus、Gretchen Helfrich、Mari Ruti、Maria Spiropulu、Mark Trodden，当然还有我的家人。（这些智慧常以“你当时在想什么？”的形式出现。）最后，我要感谢我 GR 课堂上的学生们——本书所采用的种种策略首先在他们身上接受检验——并感谢我的学生与合作者们，原谅我为了写书而在本该做研究时的心不在焉。

自己写过教科书的朋友们告诉我，书的第一版有时会包含错误。万一这种不太可能发生的事情真的在此发生，勘误表将保存在本书网站：\texttt{http://spacetimeandgeometry.net/}。该网站还将包含读者可能感兴趣的其他链接。

在写作本书期间，我得到了 National Science Foundation、Department of Energy、Alfred P.~Sloan Foundation 以及 David and Lucile Packard Foundation 的支持。

\begin{flushright}
Sean Carroll\\
芝加哥，伊利诺伊\\
2003 年 6 月
\end{flushright}
